\section{Edge Deployment Architecture and Testbed}
\label{sec:architecture}

\subsection{Two-Stage Asynchronous Edge Architecture}

To operate effectively within constrained edge budgets, our intrusion detection pipeline adopts a tiered, two-stage architecture distributed across NVIDIA Jetson Orin Nano Super boards, illustrated in Fig.~\ref{fig:architecture}. 

\begin{figure*}[t]
  \centering
  \resizebox{0.92\textwidth}{!}{%
  \begin{tikzpicture}[
      font=\small, >=Stealth, every node/.style={align=center},
      pc/.style ={draw=gray!60, thick, rounded corners, fill=gray!8, minimum width=11.6cm, minimum height=1.1cm, inner sep=3pt},
      jet/.style={draw=accGreen!80!black, thick, rounded corners, fill=accGreen!8, text width=4.6cm, minimum height=1.8cm, inner sep=4pt},
      eng/.style={draw=accBlue, thick, rounded corners, fill=accBlue!5, text width=2.1cm, minimum height=0.9cm, inner sep=2pt, font=\footnotesize},
      lbl/.style={font=\footnotesize, fill=white, inner sep=1.5pt, align=center, text=accBlue!75!black},
      arr/.style ={-{Stealth[length=2.5mm]}, thick, draw=gray!60!black},
  ]
    \node[pc]  (pc)  at (0, 0.0)   {\textbf{Network Traffic Generator / Ingestion Gateway}\\[2pt]
        Streaming Flow Replayer $\rightarrow$ ZeroMQ Publisher (\texttt{tcp://:5555})};
    \node[jet] (j1)  at (-4.5,-2.8) {\textbf{Jetson Orin Nano Super \#1}\\[2pt] 
        \textbf{Stage 1: Anomaly Gate}\\[2pt]
        scikit-learn Autoencoder + Scaler\\ (In-memory streaming filter)};
    \node[jet] (j2)  at ( 4.5,-2.8) {\textbf{Jetson Orin Nano Super \#2}\\[2pt] 
        \textbf{Stage 2: Classifier Node}\\[2pt]
        Feature-matrix builder $\rightarrow$\\ Swappable Inference Engines};
    \node[eng] (onnx)  at (3.3,-5.1) {\textbf{ONNX Runtime}\\ (\texttt{.onnx} float32)};
    \node[eng] (numpy) at (5.7,-5.1) {\textbf{NumPy Engine}\\ (\texttt{.npz} float32/64)};
    
    \draw[arr] (j2.south -| onnx)  -- (onnx.north);
    \draw[arr] (j2.south -| numpy) -- (numpy.north);
    \node[lbl] at (4.5,-6.0) {Dynamic switch: \texttt{--engine onnx\,|\,numpy}};

    \draw[arr] (pc.south -| j1) -- node[left=3pt, font=\footnotesize, text=accBlue!75!black]{\texttt{flow-stream}} (j1.north);
    \draw[arr] (j1.east) -- node[above=2pt, font=\footnotesize, text=accBlue!75!black]{ZeroMQ PUSH/PULL}
                            node[below=2pt, font=\scriptsize, text=gray!60!black]{(Zero-copy NumPy buffer: \texttt{copy=False})} (j2.west);
  \end{tikzpicture}}
  \caption{Two-stage distributed edge architecture: Anomaly Gate deployed on Jetson~\#1 filters benign background traffic, while the heavy tree-ensemble classifier operates on Jetson~\#2 with runtime-interchangeable inference backends linked via brokerless ZeroMQ.}
  \label{fig:architecture}
\end{figure*}

\begin{enumerate}
    \item \textbf{Stage 1 (Anomaly Gate on Jetson~\#1):} Incoming flows are ingested from the ZeroMQ stream endpoint \texttt{flow-\allowbreak stream}. An autoencoder reconstructs the flow vector; if the mean squared reconstruction error falls below a pre-calibrated threshold $\tau_{\text{gate}}$, the record is deemed benign and dropped immediately without incurring classifier cost. Only suspicious flows exceeding $\tau_{\text{gate}}$ are forwarded via a brokerless ZeroMQ socket directly to Jetson~\#2 using zero-copy contiguous NumPy array buffer frames (\texttt{copy=False}).
    \item \textbf{Stage 2 (Classifier Node on Jetson~\#2):} Flows forwarded by Stage 1 are ingested by the secondary board directly into contiguous memory buffers without intermediate string serialization, serving a 200-tree random forest ensemble. Crucially, Jetson~\#2 exposes an in-process runtime contract: the system can switch seamlessly between ONNX Runtime (C++ SIMD backend) and our native NumPy engine without modifying upstream transport or downstream alerting.
\end{enumerate}

\subsection{Decoupling In-Process Serving from Network Ingestion Harness}

A frequent pitfall in empirical edge evaluations is conflating the performance of the machine learning inference engine with network transport serialization overheads. In our evaluation, we deliberately separate the \textbf{In-Process Inference Engine} from the \textbf{Network Ingestion Harness}:
\begin{itemize}
    \item \textbf{In-Process Inference Benchmark:} Directly evaluates the compute kernels and memory subsystem on the Jetson Orin Nano Super boards under isolated memory buffers. Under this setup, the autoencoder gate achieves $21{,}191\,\text{flows/s}$ ($0.057\,\text{ms/flow}$), while the ONNX Runtime classifier reaches $15{,}174\,\text{flows/s}$ ($0.085\,\text{ms/flow}$).
    \item \textbf{Network Ingestion Harness:} Implemented via a streaming client (\texttt{csv\_producer.py}) communicating across an ad-hoc 802.11ac Wi-Fi channel. This harness replays recorded packet traces at variable rates to evaluate system behavior under real-world network jitter.
\end{itemize}

Under maximum replay saturation using a naive Python socket loop with JSON string parsing, the streaming harness saturates at approximately $965\,\text{flows/s}$ due to Python's Global Interpreter Lock (GIL) and string allocation overhead. By contrast, adopting brokerless ZeroMQ (\texttt{pyzmq}) with direct zero-copy buffer passing (\texttt{copy=False}) eliminates serialization entirely, allowing edge nodes to communicate at hardware wire speed and seamlessly match the $21{,}191\,\text{flows/s}$ capacity proven by our in-process serving kernels.
